\subsection{Expectation of Binomial Distribution}
\label{ssec:expectation-of-binomial-distribution}

Given a binomial random variable \(X \sim \text{Binomial}(n,
p)\) with probability mass function:
\[
  P_X(r) = P(X = r) = \binom{n}{r} q^{n-r} p^{r}, \quad r =
  0, 1, 2, \ldots, n,
\]
where \( q = 1 - p \).

\bigskip

\begin{definition}
  The expectation \(E[X]\) is defined as:
  \[
    E[X] = \sum_{r=0}^n r \, P_X(r) = \sum_{r=0}^n r
    \binom{n}{r} q^{n-r} p^r.
  \]
\end{definition}

\bigskip

\begin{note}
  The term for \(r=0\) is zero since it is multiplied by
  \(r\), so we can start the sum from \(r=1\):
  \[
    E[X] = \sum_{r=1}^n r \binom{n}{r} q^{n-r} p^r.
  \]
\end{note}

\bigskip

Rewrite \(r \binom{n}{r}\) using the identity:
\[
  r \binom{n}{r} = n \binom{n-1}{r-1}.
\]

Therefore,
\[
  E[X] = \sum_{r=1}^n n \binom{n-1}{r-1} q^{n-r} p^r = n \sum_{r=1}^n \binom{n-1}{r-1} q^{n-r} p^r.
\]

\bigskip

Change the index of summation by letting \(k = r - 1\). When
\(r=1\), \(k=0\), and when \(r=n\), \(k=n-1\). So,
\[
  E[X] = n \sum_{k=0}^{n-1} \binom{n-1}{k} q^{n - (k+1)}
  p^{k+1} = n p \sum_{k=0}^{n-1} \binom{n-1}{k} q^{(n-1)-k}
  p^k.
\]

\bigskip

Notice that the sum is the binomial expansion of \((p +
q)^{n-1}\):
\[
  \sum_{k=0}^{n-1} \binom{n-1}{k} p^k q^{(n-1)-k} = (p + q)^{n-1} = 1^{n-1} = 1.
\]

\bigskip

Hence,
\[
  E[X] = n p \cdot 1 = n p.
\]

\bigskip

\begin{remark}
  The expectation of a binomial random variable \(X \sim \text{Binomial}(n,p)\) is
  \begin{equation}
    \label{eqn:expectation-of-binomial-distribution}
    \begin{displaystyle}
      \boxed{E[X] = n p.}
    \end{displaystyle}
  \end{equation}
\end{remark}

\subsection{Variance of Binomial Distribution}
\label{ssec:variance-of-binomial-distribution}

Recall that the
\hyperref[sec:variance-of-a-random-variable]{variance of a
random variable} \(X\) is defined as:
\[
  \operatorname{Var}(X) = E[X^2] - (E[X])^2.
\]

We have already derived that,
\footnote{See \autoref{eqn:expectation-of-binomial-distribution}}
\[
  E[X] = n p.
\]

\bigskip

Next, we compute \(E[X^2]\):
\[
  E[X^2] = \sum_{r=0}^n r^2 P_X(r) = \sum_{r=0}^n r^2 \binom{n}{r} q^{n-r} p^r,
\]
where \(q = 1-p\).

\bigskip

Rewrite \(r^2\) as \(r(r-1) + r\), so that
\begin{align*}
  E[X^2] 
&= \sum_{r=0}^n \bigl[r(r-1) + r\bigr] \binom{n}{r} q^{n-r} p^r \\
&= \sum_{r=0}^n r(r-1) \binom{n}{r} q^{n-r} p^r + \sum_{r=0}^n r \binom{n}{r} q^{n-r} p^r.
\end{align*}

\bigskip

We already know the second sum is \(E[X] = n p\). Focus on the first sum:
\[
  S = \sum_{r=0}^n r(r-1) \binom{n}{r} q^{n-r} p^r.
\]

\bigskip

Using the combinatorial identity:
\[
  r(r-1) \binom{n}{r} = n (n-1) \binom{n-2}{r-2},
\]
we get
\[
  S = n (n-1) \sum_{r=2}^n \binom{n-2}{r-2} q^{n-r} p^r.
\]

\bigskip

Change the index by letting \(k = r - 2\). When \(r=2\), \(k=0\), and when \(r=n\), \(k = n-2\):
\begin{align*}
  S 
&= n (n-1) \sum_{k=0}^{n-2} \binom{n-2}{k} q^{n - (k+2)} p^{k+2} \\
&= n (n-1) p^2 \sum_{k=0}^{n-2} \binom{n-2}{k} q^{(n-2)-k} p^k.
\end{align*}

\bigskip

Recognize the sum as the binomial expansion of \((p + q)^{n-2} = 1^{n-2} = 1\):
\[
  S = n (n-1) p^2.
\]

\bigskip

Putting it all together:
\begin{align*}
  E[X^2] &= S + E[X] \\
         &= n (n-1) p^2 + n p.
\end{align*}

\bigskip

Finally, compute the variance:
\begin{align*}
  \operatorname{Var}(X) 
&= E[X^2] - (E[X])^2 \\
&= n (n-1) p^2 + n p - (n p)^2 \\
&= n (n-1) p^2 + n p - n^2 p^2.
\end{align*}

\bigskip

Simplify:
\begin{align*}
  \operatorname{Var}(X) 
&= n p + n (n-1) p^2 - n^2 p^2 \\
&= n p + n^2 p^2 - n p^2 - n^2 p^2 \\
&= n p - n p^2 \\
&= n p (1 - p) \\
&= n p q.
\end{align*}

\bigskip

\begin{remark}
  The variance of a binomial random variable \(X \sim \text{Binomial}(n,p)\) is
  \[
    \boxed{
      \operatorname{Var}(X) = n p q,
    }
  \]
  where \(q = 1 - p\).
\end{remark}
